\documentclass[11pt]{article}
\usepackage[utf8]{inputenc}
\usepackage[T1]{fontenc}
\usepackage[brazil]{babel}
\usepackage{lmodern}
\usepackage{amsmath,amssymb}
\usepackage{hyperref}
\usepackage{mathptmx}

% CONFIGURAÇÃO DE TEOREMAS
\usepackage[thmmarks]{ntheorem}

\theoremstyle{plain}
\theoremheaderfont{\normalfont\bfseries}
\theorembodyfont{\normalfont}
\theoremseparator{:}
\theorempreskip{\topsep}
\theorempostskip{\topsep}

\theoremsymbol{\ensuremath{\bigcirc}}
\newtheorem{definicao}{Definição}[section]

\theoremsymbol{\ensuremath{\triangle}}
\newtheorem{nota}[definicao]{Nota}

\theoremsymbol{}
\newtheorem{proposicao}[definicao]{Proposição}

\theoremsymbol{\ensuremath{\square}}
\newtheorem{proplivre}[definicao]{Proposição}

\theoremstyle{nonumberplain}
\theoremheaderfont{\normalfont\bfseries}
\theorembodyfont{\normalfont}
\theoremsymbol{\ensuremath{\blacksquare}}
\newtheorem{demonstracao}[definicao]{Demonstração}

% Macros
\newcommand{\naturais}{\mathbb{N}}
\newcommand{\set}[1]{\{#1\}}
\newcommand{\ainf}{A^*}
\newcommand{\funcset}[1]{#1^{#1}}
\newcommand{\powerset}[1]{\mathcal{P}(#1)}

\title{Autômatos}
\author{Gian Carlo}
\date{\today}

\begin{document}
\maketitle

\section{Introdução}

Este texto apresenta meus esforços de procurar uma definição simples, esqueletal, 
e útil de autômatos. Não afirmo que tenha conseguido a utilidade, mas acho que
consigo mostrar que a definição é simples e esqueletal, no sentido de capturar
a essência dos autômatos. Para uma literatura séria, contemplando várias 
contribuições, temos os livros de Eilenberg, de Hopcroft e de Jean Éric Pin.

O estilo do texto é sem estilo. As vezes alterno a pessoa da linguagem, entre a primeira singular e plural. As vezes busco ser impessoal. 

\section{Investigação}

Para um $n \in \naturais$, definimos \begin{equation}
    [n] = \set{1, 2, ..., n}.
\end{equation}
Com esta notação, podemos considerar um conjunto $A \neq \emptyset$ dado por \begin{equation}
    A = \set{a_i : i \in [n]}
\end{equation}
que chamaremos de alfabeto. Por fim, obtemos um conjunto de \textbf{palavras} sobre $A$
dado por \begin{equation}
    \ainf = \set \varepsilon \cup \bigcup_{n \in \naturais} A^n
\end{equation}
onde $A^n = A \times A \times ... \times A$ (produto cartesiano $n$ vezes) 
e $\varepsilon$ é a palavra vazia.

Definindo uma operação sobre $\ainf$ dada por \begin{equation}
    \begin{array}{cccc}
        \_ : & \ainf \times \ainf & \to & \ainf \\
        & (x, y) & \mapsto & xy
    \end{array}
\end{equation}
obedecendo $x\varepsilon = \varepsilon x = x$ para todo $x \in \ainf$,
obtemos um monóide livre $(\ainf, \_)$ sobre $A$. Deixemos este monóide quieto
por enquanto.

Denotaremos por $\funcset{[s]}$ o conjunto de todas as funções de $[s]$ em $[s]$. Nosso autômato será uma função \begin{equation}
    \begin{array}{cccc}
        \phi : & A & \to & \funcset{[s]} \\
         & a & \mapsto & f_a
    \end{array}
\end{equation}
sem maiores restrições. Esta definição não é novidade, já estando presente na literatura.

Não temos noção ainda de estado inicial. Para isto, iremos introduzir a seguinte notação.
Sejam $x \in [p]$ e $a \in A$. Definimos \begin{equation}
    \phi(x \mid a) = f_a(x)
\end{equation}
obedecendo \begin{equation}
    \phi(x \mid ab) = \phi(\phi(x \mid a) \mid b).
\end{equation}
Assim, $x$ cumpre o papel de um estado inicial. Ainda precisamos de estados finais.
Neste caso, apenas exigiremos que $\emptyset \neq F \subseteq [p]$ dado seja um 
conjunto de estados finais. Finalmente, podemos definir a linguagem de um autômato
como \begin{equation}
    L(\phi, x, F) = \set{w \in \ainf : \phi(x \mid w) \in F}.
\end{equation}

Estando diante de uma definição de autômato e linguagem, podemos fazer perguntas mais interessantes. Primeiro, uma vez carregada as dependências de estados inicial e finais para a linguagem, agora podemos pensar em formas de descrever autômatos ou estruturas tais que, variando ou fixando estas dependências, obtemos resultados iguais ou diferentes.

Por exemplo, poderíamos nos perguntar o que acontece se fixarmos os estados finais e variamos o estado inicial do autômato. Numa equação, estaríamos estudando o comportamento de uma função $v$ no conjunto das partes de $\ainf$ ($\powerset{\ainf}$) dada por \begin{equation}
    \begin{array}{cccc}
        v : & [s] & \to & \powerset{\ainf} \\
         & n & \mapsto & L(\phi, n, F)
    \end{array}
\end{equation}
fixando $\phi$ e $F$. Há autômatos cujo o comportamento da linguagem, sobre esta ótica, é previsível? Claramente, quando $n \in F$, teremos $\varepsilon \in L(\phi, n, F)$. Mas e as demais palavras? Que tipos de simetrias e assimetrias entram em jogo?

Como $\ainf$ é enumerável, temos $\powerset{\ainf}$ não-enumerável. De imediato pensamos em números reais ou complexos. Não é novidade que estruturas nestes espaços elucidem propriedades para fenômenos discretos, então há aí também mais uma superfície de estudo?

Uma questão inicial bem simples de se fazer seria: Se sabemos que variando $n \in [s]$, todas as nossas linguagens são translações uma das outras por palavras fixas, conseguimos determinar algo sobre $F$ e $\phi$? Se $s = 3$, por exemplo, e digamos que \begin{equation}
    L(\phi, 1, F) = aL(\phi, 2, F) = bL(\phi, 3, F),
\end{equation}
onde $xL(\phi, n, F) = \set { xp : p \in L(\phi, n, F)}$, o que isso determina de $\phi$ e $F$? Com este exemplo pequeno, seríamos capazes de reconstruir parte de $\phi$ e $F$ relativos à equação por inspeção e estudo de casos, mas e num caso geral?

\section{Síntese}

\begin{nota} \label{nota:notacao_intervalo_discreto}
    Seja $n \in \naturais$. Denotaremos por $[n]$ o conjunto $\set{1, 2, ..., n}$.
\end{nota}

\begin{definicao}[Autômato] \label{def:automato}
    Seja $s \in \naturais$ e $A \neq \emptyset$ um conjunto. Um \textbf{autômato} é uma função \begin{equation}
        \begin{array}{cccc}
            \phi : & A & \to & \funcset{[s]} \\
            & a & \mapsto & f_a
        \end{array}
    \end{equation}
    onde $f_a$ é uma função $[s] \to [s]$.
\end{definicao}

\begin{nota} \label{nota:notacao_estado_inicial}
    Introduziremos mais algumas notações: \begin{enumerate}
        \item Dados $x \in [s]$ e $a \in A$, faremos \begin{equation}
            \phi(x \mid a) = f_a(x)
        \end{equation}
        Note que $\phi(x \mid a) \in [s]$. Ou seja, podemos expandir a notação por meio de uma recursão.
        \item Dado $x \in [s]$ e $a, b \in A$, faremos \begin{equation}
            \phi(x \mid ab) = \phi(\phi(x \mid a) \mid b)
        \end{equation}
        Que de fato faz sentido, pois $\phi(x \mid a) \in [s]$ e assim podemos substituir o $x$ em $\phi(x \mid b)$.
    \end{enumerate}
\end{nota}

\begin{proplivre} \label{prop:composicao_automato}
    Seja $\phi$ um autômato, $x \in [s]$ e $a, b \in A$. Temos que \begin{equation}
        \phi(x \mid ab) = (\phi(b) \circ \phi(a))(x).
    \end{equation}
\end{proplivre}

\begin{nota}
    Denotaremos por $\ker \phi$ o conjunto \begin{equation}
        \ker \phi = \set{a \in A : \phi(a) = id},
    \end{equation}
    por $A \leq B$ o fato de que $A$ é um subgrupo de $B$, e por $S_s$ o grupo de permutações (funções bijetoras) de $[s] \to [s]$. Imediatamente temos que $S_s \subseteq \funcset{[s]}$.
\end{nota}

\begin{proposicao}[Autômato Subgrupo] \label{prop:automato_subgrupo}
    Seja $\phi$ um autômato e $[s]$ seu conjunto de estados. Se $\phi(A) \leq S_s$, então valem: \begin{enumerate}
        \item Para todo $n \in [s]$, existem $a, b \in A$ tais que $\phi(n | ab) = n$;
        \item Para todo $a \in \ker \phi$ e $n \in [s]$, temos $\phi(n | a) = n$; e
        \item Para todo $n \in [s]$ e $a, b \in A$, existe $c \in A$ tal que $\phi(n | ab) = \phi(n | c)$.
    \end{enumerate}
\end{proposicao}

\begin{demonstracao}
    Para (1), tome $a \in A$ e denote por $\sigma_a$ a permutação $\phi(a)$. Como $\phi(A) \leq S_s$, temos que $\sigma_a^{-1} \in \phi(A)$, ou seja, existe $b \in A$ tal que $\phi(b) = \sigma_a^{-1}$. Assim, para todo $n \in [s]$, temos \begin{equation}
        \phi(n | ab) = (\phi(b) \circ \phi(a))(n) = (\sigma_a^{-1} \circ \sigma_a)(n) = n.
    \end{equation}
    
    Para (2), seja $a \in \ker \phi$ e $n \in [s]$. Então $\phi(a) = id$, e assim \begin{equation}
        \phi(n | a) = \phi(a)(n) = id(n) = n.
    \end{equation}
    
    Para (3), tome $n \in [s]$ e $a, b \in A$. Denote por $\sigma_a$ e $\sigma_b$ as permutações $\phi(a)$ e $\phi(b)$, respectivamente. Como $\phi(A) \leq S_s$, temos que $\sigma_a \circ \sigma_b \in \phi(A)$, ou seja, existe $c \in A$ tal que $\phi(c) = \sigma_a \circ \sigma_b$. Assim, temos \begin{equation}
        \phi(n | ab) = (\phi(b) \circ \phi(a))(n) = (\sigma_a \circ \sigma_b)(n) = \phi(c)(n) = \phi(n | c).
    \end{equation}
\end{demonstracao}

\begin{nota}
    Introduziremos terminologia para as propriedades da proposição anterior, respectivamente: \begin{enumerate}
        \item Diremos que o autômato em questão sempre tem \textbf{transição inversível};
        \item Todo $a \in \ker \phi$ é chamado de \textbf{barulho}; 
        \item Diremos que o autômato é \textbf{transitivo}.
    \end{enumerate}
\end{nota}

\begin{nota}
    Seja $A$ um conjunto. Denotaremos por $A^p$ o produto cartesiano \begin{equation}
        A \times A^{p-1}
    \end{equation}
    onde $A^1 = A$. Por extenso, temos que $A^p$ é o produto cartesiano de $p$ cópias de $A$: \begin{equation}
        A^p = \underbrace{A \times A \times ... \times A}_{p}
    \end{equation}
\end{nota}

\begin{definicao}[Conjunto de Palavras]
    Denotaremos por $\ainf$ o conjunto \begin{equation}
        \ainf = \set \varepsilon \cup \bigcup_{p = 1}^{\infty} A^p
    \end{equation}
    que chamaremos de \textbf{palavras de $A$}. Neste contexto, $A$ será chamado de conjunto de \textbf{símbolos}. Aqui $\varepsilon$ representa uma \textbf{palavra vazia}, sem símbolos.
\end{definicao}

\begin{nota}
    Se definimos uma operação de concatenação em $\ainf$, que denotaremos por $+$, e exigimos que \begin{equation}
        p + \varepsilon = \varepsilon + p = p
    \end{equation}
    para todo $p \in \ainf$, obtemos o que na literatuara se chama de \textbf{monóide livre} sobre $A$, dado pela tupla $(\ainf, +)$.
\end{nota}

\begin{nota}
    Sejam $\phi : A \to \funcset{[s]}$ um autômato, $i \in [s]$, e $p \in \ainf$. Escrevendo $p$ por extenso como \begin{equation}
        p = (p_1, p_2, ..., p_k)
    \end{equation}
    extendemos nossa definição de $\phi(x | y)$ fazendo \begin{equation}
        \phi(i | p) = \phi(i | p_1p_2...p_k).
    \end{equation}
\end{nota}

\begin{definicao}[Linguagem] \label{def:linguagem}
    Seja $\phi$ um autômato, $i \in [s]$ e $\emptyset \neq F \subseteq [s]$. A \textbf{linguagem} de $(\phi, i, F)$, que denotaremos por $L(\phi, i, F)$, é o conjunto \begin{equation}
        L(\phi, i, F) = \set{p \in \ainf : \phi(i \mid p) \in F}.
    \end{equation}
\end{definicao}

\section{Conclusões Provisórias}

A ideia no momento é ir alternando entre as investigações e a síntese. Não sei se irei modificar a parte da investigação para mencionar o que é explicitado na síntese, ou se deixo como foi escrito para fins de registro histórico. Claramente o estudo ainda está extremamente superficial.

Estou ciente de que nada aqui é novidade. Este tipo de trabalho serve como motivação para mim mesmo. Sinto que aprendo desta forma, buscando eu mesmo trilhar estes caminhos.

\end{document}
